\documentclass[a4paper,12pt]{letter}

\usepackage{amsmath}
\usepackage{amssymb}
\usepackage{nopageno}
\usepackage{amsmath}
\usepackage{enumitem}
% \usepackage{mathtools}
\usepackage{eqnarray,amsmath}
\usepackage[a4paper, total={7.5in, 11.25in}]{geometry}

\begin{document}

Name: \hrulefill\ Neptun code: \hrulefill

\bigskip

\begin{center}
(KTXFI2EBNF) Physics II. Exam~1 \\
December~16,~2025 
\end{center}

\bigskip
 
\begin{enumerate}[label=\Roman*.]
   
\item Moving Reference Frames. Special Theory of Relativity

\begin{enumerate}[label=\arabic*.]
\setcounter{enumii}{0} % Számláló beállítása

%\item State the Galilean principle of relativity. \hfill (5~points)

\item  Write down the Lorentz transformation equations used in the special theory of relativity. Specify the exact physical meaning of the symbols used in the formulas, including the reference frame in which they are defined. \hfill (5~points)
  
\item Write down Einstein’s relation expressing mass–energy equivalence. What is its physical meaning? \hfill (5~points)
 
\end{enumerate}

% ------------------------------------------------------------------------------------------------------------------------------------

\item Limits of the Classical Conceptual Framework

\begin{enumerate}[label=\arabic*.]
\setcounter{enumii}{2} % Számláló beállítása
  
\item What is the photoelectric effect? Explain the mechanism of the photoelectric effect. \hfill (5~points)
  
\item Describe the Compton experiment and provide a qualitative but precise explanation of the phenomenon of Compton scattering. \hfill (5~points)

\item Formulate the de Broglie principle in your own words. \hfill (5~points)

%\item What is meant by the dual nature of electromagnetic radiation? Describe the meaning of this expression. \hfill (5~points)
 
\end{enumerate}
  
% ------------------------------------------------------------------------------------------------------------------------------------

\item Elements of Quantum Mechanics. Physics of Condensed Matter. Theory of Electrical Conduction

\begin{enumerate}[label=\arabic*.]
\setcounter{enumii}{5} % Számláló beállítása
  
\item What does the Heisenberg uncertainty principle express? To which types of physical quantity pairs does it apply? \hfill (5~points)

%\item Describe and explain the electrical conduction mechanism of semiconductors using the band theory of solids. What characterizes semiconductors from the point of view of electrical conduction? \hfill (5~points)

\item What is the Peltier effect and what is its explanation? Draw a diagram. Interpret the effect and provide the physical explanation of the phenomenon. \hfill (5~points)

\end{enumerate}

% ------------------------------------------------------------------------------------------------------------------------------------
  
\item Development of Atomic Theory

\begin{enumerate}[label=\arabic*.]
\setcounter{enumii}{7} % Számláló beállítása
  
\item  State the Pauli exclusion principle. \hfill (5~points)

\item Give the physical interpretation of the principal quantum number. What is the principal quantum number? What does it represent? What is its physical meaning? What values can it take? \hfill (5~points)

\item Given the element with atomic number 10 in the periodic table, write down its electronic configuration using the standard atomic physics notation. \hfill (5 points)

\end{enumerate}

\end{enumerate}

% ------------------------------------------------------------------------------------------------------------------------------------

Grades: 0–25~points: Fail (1), 26–30 points: Pass (2), 31–37 points: Satisfactory (3), 38–45~points: Good (4), 46–50 points: Excellent (5).

\rightline{Dr.~Gabor FACSKO}
\rightline{\textit{facsko.gabor@uni-obuda.hu}}

\vfill

\end{document}
